% -------------------------------------------------------
%  English Abstract
% -------------------------------------------------------

\begin{latin}

\begin{center}
\textbf{Abstract}
\end{center}
\baselineskip=.8\baselineskip

In non-contrast CT, soft-tissue contrast is low, and the performance of segmentation models, particularly for vessels, drops compared with contrast-enhanced images. The scarcity of annotated data for this phase and the difficulty of registering the different scan phases, caused by anatomical changes and the time gap between acquisitions, further complicate the problem.

In this work, an anatomy-aware preprocessing for registering the scan phases was first proposed, which reduced the number of large registration errors (organ-centroid displacement above \عدد{10} mm) for the most accurate method examined, \lr{DEEDS}, by \عدد{65}\%. The registered pairs were then used to train the proposed model, \lr{PLC-Net}. The model takes a non-contrast image together with the target scan phase and the target enhancement level of the organs in Hounsfield units, and generates the same patient's image in the target phase. Experiments were conducted on the \lr{VinDr-Multiphase} dataset with a patient-level split of \عدد{97}, \عدد{20} and \عدد{20} cases for training, validation and testing, and the outputs were evaluated with image metrics and with automatic segmentation as the downstream task.

The output of \lr{PLC-Net} increased the mean \lr{Dice} of portal-vein segmentation from \عدد{0.55} on the non-contrast image to \عدد{0.72}. Among the image metrics, the organ-level median Hounsfield-unit error showed the strongest correlation with the segmentation improvement. With an error of \عدد{14.55} HU on this metric, \lr{PLC-Net} performed best among the compared methods, while the closest method, \lr{ResViT}, reached \عدد{15.17} HU with about \عدد{9} times more parameters. These results show that the proposed approach, using a lightweight conditional model, can generate contrast-enhanced images that are useful for the downstream task, and offers a way to use non-contrast images for abdominal organ segmentation.

\bigskip\noindent\textbf{Keywords}:
Multi-organ segmentation, Abdominal CT, Non-contrast imaging, Inter-phase
registration, Conditional image synthesis, Downstream evaluation

\end{latin}
